% Generated draft: requires mathematical and source review.
\documentclass[11pt]{article}
\usepackage[margin=25mm]{geometry}
\usepackage{fontspec}
\setmainfont{Latin Modern Roman}
\usepackage{amsmath,amssymb,mathtools,unicode-math}
\setmathfont{Latin Modern Math}
\usepackage{enumitem,needspace,xurl,hyperref}
\setlength{\parindent}{0pt}
\setlength{\parskip}{5pt}
\setlength{\emergencystretch}{3em}
\newcommand{\N}{\mathbb{N}}
\newcommand{\Z}{\mathbb{Z}}
\newcommand{\Q}{\mathbb{Q}}
\newcommand{\R}{\mathbb{R}}
\newcommand{\C}{\mathbb{C}}
\newcommand{\F}{\mathbb{F}}
\newcommand{\A}{\mathbb{A}}
\newcommand{\PP}{\mathbb{P}}
\newcommand{\E}{\mathbb{E}}
\newcommand{\eps}{\varepsilon}
\newcommand{\dd}{\,\mathrm{d}}
\begin{document}
\section*{UnsolvedMath ID: 800814}
\textbf{Source problem code:} catalogue entry
\subsection{Definitions and mathematical statement}
The supplied Problem 13 asks for the construction of anisotropic function spaces adapted to ideal hydrodynamics, with greater smoothness along flow lines than in transverse directions, and for existence and uniqueness theorems in those spaces in any spatial dimension.

For terminology, if \(u(t,x)\) is a sufficiently regular velocity field on a spatial domain or manifold \(M\), its flow trajectories satisfy
\[
\frac{\dd}{\dd t}\Phi(t,a)=u(t,\Phi(t,a)),\qquad \Phi(0,a)=a.
\]
For a scalar function \(f\), spatial differentiation in the instantaneous flow direction is \(u(t,\cdot)\cdot\nabla f\). These notions describe the requested directional regularity but do not define a function space. In particular, differentiation along a space--time trajectory, \(\partial_t+u\cdot\nabla\), differs from spatial differentiation along an instantaneous streamline; the supplied statement does not specify which interpretation is intended.

\textbf{Research problem.} Construct anisotropic function spaces expressing greater regularity in the intended flow direction than in transverse directions. Formulate and prove existence and uniqueness results for ideal hydrodynamics in these spaces, in every spatial dimension covered by the intended formulation. The construction must specify the spaces, their regularity parameters, how their directional structure is determined by the velocity, the admissible initial data, and the solution and uniqueness classes.

\textbf{Unresolved qualifications.} The supplied record does not specify the hydrodynamic equations, the spatial domain or manifold, boundary conditions, the meaning of directional smoothness, or whether existence is local or global in time. It also does not specify the admissible dimensions or how flow directions are treated at zeros of the velocity. Consequently, it supplies a research program rather than a fully quantified existence and uniqueness assertion. A precise theorem cannot be extracted without additional hypotheses; no particular equation, domain, regularity threshold, or time interval is imposed here.
\subsection{Short English statement}
Construct function spaces that allow greater smoothness along fluid flow lines than across them, and establish existence and uniqueness for ideal hydrodynamics in those spaces in any dimension. The source leaves the equations, spaces, geometry, boundary conditions, and time horizon unspecified, so a fully precise theorem remains to be formulated.
\subsection{Sources}
\begin{itemize}
\item Local UnsolvedMath registry, permanent ID \(800814\), source code \texttt{catalogue entry}, title: \emph{Geometric Hydrodynamics in Open Problems --- Problem 13}.
\item Khesin, Misiolek, Modin, \emph{Geometric Hydrodynamics in Open Problems} (2022), Problem 13. Supplied source URL: \url{https://arxiv.org/abs/2205.01143}. The record attributes the problem to these authors and gives the proposed year as 2022.
\item Supplied provenance: accessed 2026-10-16; extraction method \texttt{arxiv\_latex}; upstream revision \texttt{d371ee917ac9647fbb06cd6fc529deec6f3dda82}. These are supplied provenance details, not an independent source-access claim.
\item The local record labels the problem \emph{open} and gives an update timestamp of \texttt{2026-10-16T14:00:00Z}. This is supplied dated curation, not an independently verified assessment of the literature or mathematical status.
\end{itemize}
\end{document}
